\documentclass[10pt,a4paper]{amsart}
\usepackage[margin=19mm]{geometry}
\usepackage{amsmath,amssymb,mathtools}
\usepackage[scr=rsfs,cal=euler]{mathalfa}
\usepackage{xfrac}
\usepackage[unicode,hidelinks,bookmarks=false]{hyperref}
\newcommand{\ie}{i.e.\ }
\newcommand{\cf}{cf.\ }
\newcommand{\resp}{respectively\ }
\newcommand{\Subsection}{\subsection}
\providecommand{\nobreakdash}{\nobreak\hbox{-}}
\setlength{\emergencystretch}{2em}
\multlinegap0pt
\usepackage{mathtools,mathrsfs}
\usepackage{mathptmx}
\usepackage{enumitem,booktabs,array}
\setlength{\emergencystretch}{2em}
\numberwithin{equation}{section}
\newtheorem{definition}{Definition}[section]
\newtheorem{example}[definition]{Example}
\newtheorem{theorem}[definition]{Theorem}
\newtheorem{proposition}[definition]{Proposition}
\newtheorem{lemma}[definition]{Lemma}
\newtheorem{corollary}[definition]{Corollary}
\newtheorem*{theoremA}{Theorem A}
\newtheorem{remark}[definition]{Remark}
\renewcommand{\proofname}{Proof}
\renewcommand{\abstractname}{Abstract}
\renewcommand{\refname}{References}
\DeclareMathOperator{\Ext}{Ext}
\DeclareMathOperator{\Hom}{Hom}
\DeclareMathOperator{\Tor}{Tor}
\DeclareMathOperator{\pd}{pd}
\DeclareMathOperator{\depth}{depth}
\DeclareMathOperator{\rad}{rad}
\newcommand{\Z}{\mathbb Z}
\newcommand{\N}{\mathbb N}
\newcommand{\D}{\operatorname{D}}
\newcommand{\gmod}[1]{#1\text{-}\mathrm{gr}}
\newcommand{\umod}[1]{#1\text{-}\mathrm{mod}}
\newcommand{\tw}[2]{{}^{#1}\!#2}
\newcommand{\Id}{\mathrm{Id}}
\hypersetup{pdftitle={Self-extensions over quantum complete intersections},
  pdfsubject={Compact manuscript; left modules and arbitrary commutation parameters}}
\begin{document}
\title[The Auslander--Reiten conjecture for quantum complete inters]{The Auslander--Reiten conjecture for quantum complete intersections}
\author{Weiheng Xia}
\date{}
\maketitle
\noindent\emph{Source and licence.} The Auslander--Reiten conjecture for quantum complete intersections. Version arXiv:2609.24007v1 (CC BY 4.0). This material is distributed under the Creative Commons Attribution 4.0 International licence, http://creativecommons.org/licenses/by/4.0/, as recorded in the arXiv OAI licence metadata for this exact version and shown on the abstract page https://arxiv.org/abs/2609.24007v1. Full rights notice: licenses/arxiv-2609.24007-CC-BY-4.0.txt.\par\smallskip
\noindent\textit{Editorial scope.} Counted unit: the original \texttt{theorem} environment \texttt{thm:main} at source lines 475--482 together with its complete proof at lines 484--589; the delivered document keeps source lines 98--590 so that every referenced label and every citation resolves inside it. Native editable source is the paper's own concatenated TeX; the AMS wrapper is editorial formatting only. No publisher class, upstream script or unrelated artwork is redistributed.
\begin{equation}\label{eq:QCI}
 A=A(\mathbf a,\mathbf q)=
 \frac{k\langle x_1,\ldots,x_c\rangle}
 {(x_i^{a_i}\ (1\leq i\leq c),\quad
   x_ix_j-q_{ij}x_jx_i\ (1\leq i<j\leq c))}.
\end{equation}
Our main result is the following projectivity criterion.

\begin{theoremA}
For every field \(k\), every algebra \(A\) in \eqref{eq:QCI}, and every
finite-dimensional left \(A\)-module \(M\),
\[
 \Ext_A^1(M,M)=0=\Ext_A^2(M,M)
 \quad\Longrightarrow\quad M\text{ is projective}.
\]
\end{theoremA}

Thus ARC and TC2 hold for \eqref{eq:QCI} with arbitrary nonzero
commutation parameters. In particular, they hold for the Liu--Schulz
family. The proof uses the gradability theorem of Amiot--Oppermann
\cite[Corollary 4.5]{AO13}, the graded twists introduced by Zhang
\cite{Zhang96}, and the commutative criterion of
\cite[Theorem 4.2]{AB00}. We construct a finite sequence of graded twists leading to a commutative truncated polynomial algebra. At each step, the first and second ordinary self-extension groups remain zero, and projectivity is preserved and reflected.

Section~\ref{sec:prelim} records the algebra and homological facts used
below. Section~\ref{sec:transport} establishes the twist and Ext formulas.
Section~\ref{sec:main} proves Theorem A, gives its consequences, and
contains a complete four-generator calculation.

\medskip\noindent\textbf{Conventions.}
Algebras and modules over \(k\) are finite-dimensional unless stated
otherwise, and all modules are left modules. Write \(\umod{B}\) for
ordinary modules and \(\gmod{B}\) for \(\Z\)-graded modules with
degree-zero morphisms; \(\Ext_B\) denotes ordinary Ext. Morphisms act on
the right and compose from left to right:
\((x)(fg)=((x)f)g\). Functors compose from right to left:
\((FG)(X)=F(G(X))\).

\section{Preliminaries}\label{sec:prelim}

\subsection{Quantum complete intersections}

For \(A\) in \eqref{eq:QCI}, set
\[
 I=\{\alpha\in\N^c:0\leq\alpha_i<a_i\},\qquad
 x^\alpha=x_1^{\alpha_1}\cdots x_c^{\alpha_c},\qquad
 \kappa(\alpha,\beta)=\prod_{u<v}q_{uv}^{-\alpha_v\beta_u}.
\]
All inequalities between exponent vectors are coordinatewise.

\begin{lemma}\label{lem:structure}
The elements \(x^\alpha\), \(\alpha\in I\), form a basis of \(A\), with
\begin{equation}\label{eq:PBWproduct}
 x^\alpha x^\beta=
 \begin{cases}
 \kappa(\alpha,\beta)x^{\alpha+\beta},&\alpha+\beta\in I,\\
 0,&\alpha+\beta\notin I.
 \end{cases}
\end{equation}
The algebra \(A\) is local and Frobenius, hence self-injective.
For each \(j\), it has the coordinate grading
\(\deg_jx_i=\delta_{ij}\). Every diagonal assignment
\((x_i)\delta_\lambda=\lambda_i x_i\), with \(\lambda_i\in k^\times\),
defines an automorphism preserving these gradings.
\end{lemma}

\begin{proof}
Reordering by \(x_vx_u=q_{uv}^{-1}x_ux_v\) shows that the displayed
monomials span and gives \eqref{eq:PBWproduct}. On a vector space with
basis indexed by \(I\), this formula defines an associative multiplication:
\[
 \kappa(\alpha,\beta)\kappa(\alpha+\beta,\gamma)
 =\kappa(\beta,\gamma)\kappa(\alpha,\beta+\gamma),
\]
If \(\alpha+\beta+\gamma\notin I\), both bracketings are zero.
The basis elements corresponding to the standard unit vectors
satisfy the defining relations of \(A\).
The resulting homomorphism from \(A\) sends the spanning monomials
to distinct basis elements, proving their linear independence.

The ideal \(J=(x_1,\ldots,x_c)\) is nilpotent and \(A/J=k\), so \(A\)
is local. Let \(\omega=(a_1-1,\ldots,a_c-1)\), and let \(\ell:A\to k\)
extract the coefficient of \(x^\omega\). Relative to the bases
\(\{x^\alpha\}\) and \(\{x^{\omega-\beta}\}\), the pairing
\((a,b)\mapsto(ab)\ell\) has diagonal matrix with nonzero entries
\(\kappa(\alpha,\omega-\alpha)\). It is associative and nondegenerate,
so \(A\) is Frobenius. Finally, the relations are homogeneous for the
coordinate gradings and are preserved by diagonal scaling.
\end{proof}

\subsection{Graded homological facts and scalar extension}

For \(U\in\gmod{B}\), use the shift \(U(d)_r=U_{r+d}\).
Thus a degree-\(d\) map \(U\to V\) is a degree-zero map \(U\to V(d)\).

\begin{lemma}\label{lem:graded-facts}
Let \(B\) be a finite-dimensional \(\Z\)-graded algebra and
\(U,V\in\gmod{B}\).
\begin{enumerate}
\item Every \(B\)-linear map is a finite sum of homogeneous maps, and
\[
 \Hom_B(U,V)=\bigoplus_{d\in\Z}\Hom_{\gmod{B}}(U,V(d)).
\]
\item The category \(\gmod{B}\) has enough projectives. Its projective
objects are precisely the objects whose underlying modules are
projective. Every object has a projective resolution with
finite-dimensional graded terms and degree-zero differentials.
\item For every \(n\geq0\),
\begin{equation}\label{eq:Ext-decomp}
 \Ext_B^n(U,V)\simeq
 \bigoplus_{d\in\Z}\Ext_{\gmod{B}}^n(U,V(d)).
\end{equation}
\end{enumerate}
\end{lemma}

\begin{proof}
The Hom decomposition and the projectivity comparison are
\cite[Theorem 1.2.6 and Proposition 1.2.15]{Hazrat16}, applied to
\(B^{\mathrm{op}}\) to obtain left-module statements. Shifted free
modules provide enough projectives; choosing finite homogeneous
generators for successive kernels gives the stated resolutions.
For (3), take such a resolution
\(\cdots\to P_1\to P_0\to U\to0\). It is also an ordinary projective
resolution. Apply (1) termwise to its Hom complex. The differential
preserves internal degree, and direct sums are exact, so taking
cohomology gives \eqref{eq:Ext-decomp}.
\end{proof}

\begin{lemma}\label{lem:base-change}
Let \(A\) be a finite-dimensional \(k\)-algebra, \(K/k\) a field
extension, and \(M,N\in\umod{A}\). Write \(A_K=K\otimes_k A\) and
\(M_K=K\otimes_kM\). Then
\begin{equation}\label{eq:base-Ext}
 K\otimes_k\Ext_A^n(M,N)\simeq\Ext_{A_K}^n(M_K,N_K)
 \qquad(n\geq0).
\end{equation}
Moreover, \(M_K\) is projective if and only if \(M\) is projective.
\end{lemma}

\begin{proof}
Take a projective resolution of \(M\) with finitely generated terms.
The degreewise Hom--scalar-extension isomorphisms and exactness of
\(K\otimes_k-\) give \eqref{eq:base-Ext}. Scalar extension preserves
projectivity. Conversely, suppose that \(M_K\) is projective.
Choose a short exact sequence
\[
 \xi:0\longrightarrow N\longrightarrow P
 \longrightarrow M\longrightarrow0
\]
with \(P\) finite free and apply \(K\otimes_k-\).
The resulting sequence splits because \(M_K\) is projective.
Thus \(1\otimes[\xi]=0\) by \eqref{eq:base-Ext}.
The map
\[
 V\longrightarrow K\otimes_k V,
 \qquad v\longmapsto1\otimes v,
\]
is injective for every \(k\)-vector space \(V\).
Hence \([\xi]=0\), so \(M\) is projective.
\end{proof}

\section{Gradings, twists, and ordinary self-extensions}\label{sec:transport}

Starting with an ordinary rigid module, we choose a compatible grading, apply a left Zhang twist, and then forget the grading. To compare ordinary self-extension groups before and after this construction, we express them as direct sums of graded Ext groups with shifted targets. Since twisting does not in general commute with grading shifts, we first establish the precise relation between these two operations. We then use this relation to prove the vanishing result needed in the proof of the main theorem.

A module \(M\) is \emph{rigid} if \(\Ext_B^1(M,M)=0\).
The following is the gradability theorem used in each step of the
construction.

\begin{theorem}[{\cite[Corollary 4.5]{AO13}}]\label{thm:AO}
Let \(k\) be algebraically closed and let \(B\) be a finite-dimensional
\(\Z\)-graded \(k\)-algebra. Every finite-dimensional rigid ordinary
left \(B\)-module admits a compatible grading.
\end{theorem}

The left-module statement follows by applying the cited result to
\(B^{\mathrm{op}}\). No semisimplicity assumption on \(B_0\) is needed.
For an automorphism \(\alpha\) of \(B\), let \(\tw{\alpha}{M}\) denote
the ordinary left module with the same vector space and action
\begin{equation*}
 b\cdot_\alpha m=((b)\alpha)m.
\end{equation*}
If \(f:M\to N\) is \(B\)-linear, the same underlying linear map
defines a \(B\)-linear map
\(\tw{\alpha}{M}\to\tw{\alpha}{N}\).
In particular, twisting by \(\alpha\) preserves isomorphisms.

\begin{proposition}\label{prop:diagonal}
Assume that \(k\) is algebraically closed. If \(M\) is a rigid module
over a quantum complete intersection \(B\), then
\(\tw{\delta}{M}\simeq M\) as ordinary modules for every diagonal
automorphism \(\delta\), where \((x_i)\delta=\lambda_i x_i\) and
\(\lambda_i\in k^\times\).
\end{proposition}

\begin{proof}
Fix a coordinate \(i\), give \(x_i\) degree one and the other generators
degree zero, and apply Theorem~\ref{thm:AO} to grade \(M\).
For \(t\in k^\times\), the automorphism \(\delta_{i,t}\) scaling only
\(x_i\) satisfies \((b)\delta_{i,t}=t^s b\) for \(b\in B_s\).
The linear isomorphism
\[
 h_t:M\longrightarrow\tw{\delta_{i,t}}{M},
 \qquad (m)h_t=t^r m\quad(m\in M_r)
\]
is \(B\)-linear, since for \(b\in B_s\),
\[
 (bm)h_t=t^{r+s}bm
 =((b)\delta_{i,t})((m)h_t).
\]
Every diagonal automorphism is a product of these commuting coordinate
scalings. Successively twisting the resulting ordinary isomorphisms
proves the assertion. The coordinate gradings may be chosen separately;
a simultaneous \(\Z^c\)-grading is not required.
\end{proof}

For the remainder of this section the field is arbitrary, \(B\) is a
finite-dimensional \(\Z\)-graded algebra, and \(\sigma\) is a
grading-preserving automorphism. We use the following left version of
the automorphism twist of \cite{Zhang96}. On the graded vector space
\(B\), set
\begin{equation*}
 a*b=((a)\sigma^t)b,
 \qquad a\in B_s,\ b\in B_t,
\end{equation*}
and denote the resulting algebra by \(B^\sigma\).
For \(M\in\gmod{B}\), let \(F_\sigma M\) have the same graded vector
space, with action
\begin{equation*}
 b*m=((b)\sigma^r)m,
 \qquad b\in B_s,\ m\in M_r.
\end{equation*}
Both formulas extend bilinearly. On degree-zero morphisms,
\(F_\sigma\) retains the underlying linear map.

\begin{lemma}\label{lem:twist-equivalence}
These formulas define a graded algebra and an exact equivalence
\[
 F_\sigma:\gmod{B}\longrightarrow\gmod{B^\sigma}
\]
with inverse \(F_{\sigma^{-1}}\). The underlying ordinary module of
\(F_\sigma M\) is projective if and only if that of \(M\) is projective.
\end{lemma}

\begin{proof}
For \(a\in B_u\), \(b\in B_s\), and \(c\in B_t\),
\[
 (a*b)*c=((a)\sigma^{s+t})((b)\sigma^t)c=a*(b*c).
\]
The unit is unchanged, and \(\sigma\) remains a graded automorphism
of \(B^\sigma\). For \(m\in M_r\),
\[
 a*(b*m)=((a)\sigma^{r+s})((b)\sigma^r)m=(a*b)*m.
\]
If \(f:M\to N\) is a degree-zero \(B\)-homomorphism, then
\((b*m)f=((b)\sigma^r)((m)f)=b*((m)f)\), so the construction is
functorial. Twisting again by \(\sigma^{-1}\) restores the original
products and actions:
\[
 ((a)\sigma^{-t})*b=ab,
 \qquad ((b)\sigma^{-r})*m=bm.
\]
Thus \(F_{\sigma^{-1}}F_\sigma=\Id\) and
\(F_\sigma F_{\sigma^{-1}}=\Id\), with the indicated source categories.
Exactness follows because the underlying vector spaces and maps are
unchanged. The equivalence preserves and reflects graded projectivity;
Lemma~\ref{lem:graded-facts} then gives the ordinary assertion.
\end{proof}

Recall the shift convention \(M(d)_r=M_{r+d}\). Since \(\sigma\)
preserves degrees, \(\tw{\sigma^d}{M}\) retains the grading of \(M\).

\begin{lemma}\label{lem:shift-twist}
For every \(d\in\Z\), there is an identity of graded
\(B^\sigma\)-modules
\begin{equation*}
 F_\sigma\bigl(\tw{\sigma^d}{M}(d)\bigr)=F_\sigma(M)(d).
\end{equation*}
\end{lemma}

\begin{proof}
The degree-\(r\) component on either side is \(M_{r+d}\).
For \(m\in M_{r+d}\) and \(b\in B_s\), the action on the left is
\[
 b*m=((b)\sigma^r)\cdot_{\sigma^d}m
     =((b)\sigma^{r+d})m.
\]
On the right, the twist is applied before the shift, while \(m\) has
degree \(r+d\), and gives the same action.
\end{proof}

\begin{proposition}\label{prop:twisted-Ext}
For \(M\in\gmod{B}\) and \(n\geq0\), there is a natural
vector-space isomorphism
\begin{equation}\label{eq:twisted-Ext}
 \Ext_{B^\sigma}^n(F_\sigma M,F_\sigma M)
 \simeq
 \bigoplus_{d\in\Z}
 \Ext_{\gmod{B}}^n\bigl(M,\tw{\sigma^d}{M}(d)\bigr).
\end{equation}
Here the Ext group on the left is computed in the ordinary module
category.
\end{proposition}

\begin{proof}
Apply the graded decomposition of ordinary Ext from
Lemma~\ref{lem:graded-facts}, followed by
Lemma~\ref{lem:shift-twist} and the exact equivalence above:
\begin{align*}
 \Ext_{B^\sigma}^n(F_\sigma M,F_\sigma M)
 &\simeq\bigoplus_d
 \Ext_{\gmod{B^\sigma}}^n(F_\sigma M,(F_\sigma M)(d))\\
 &\simeq\bigoplus_d
 \Ext_{\gmod{B^\sigma}}^n
 \bigl(F_\sigma M,F_\sigma(\tw{\sigma^d}{M}(d))\bigr)\\
 &\simeq\bigoplus_d
 \Ext_{\gmod{B}}^n\bigl(M,\tw{\sigma^d}{M}(d)\bigr).
\end{align*}
\end{proof}

\begin{proposition}\label{prop:transport}
Let \(M\in\gmod{B}\) and \(S\subseteq\Z_{>0}\). Suppose that
\[
 \Ext_B^n(M,M)=0\quad(n\in S),
 \qquad \tw{\sigma^d}{M}\simeq M\quad(d\in\Z)
\]
as ordinary modules. Then
\[
 \Ext_{B^\sigma}^n(F_\sigma M,F_\sigma M)=0\quad(n\in S).
\]
Furthermore, \(F_\sigma M\) is projective as an ordinary
\(B^\sigma\)-module if and only if \(M\) is projective as an
ordinary \(B\)-module.
\end{proposition}

\begin{proof}
Fix \(d\in\Z\) and \(n\in S\). The assumed ordinary isomorphism
and Lemma~\ref{lem:graded-facts} give
\[
 0=\Ext_B^n(M,\tw{\sigma^d}{M})
 \simeq\bigoplus_{e\in\Z}
 \Ext_{\gmod{B}}^n\bigl(M,\tw{\sigma^d}{M}(e)\bigr).
\]
In particular, its \(e=d\) summand is zero. This holds for every
\(d\), so \eqref{eq:twisted-Ext} proves the required vanishing.
The projectivity assertion is Lemma~\ref{lem:twist-equivalence}.
\end{proof}

The isomorphisms \(\tw{\sigma^d}{M}\simeq M\) need not preserve
the chosen gradings. Together with the assumed vanishing of
\(\Ext_B^n(M,M)\), they imply that
\(\Ext_B^n(M,\tw{\sigma^d}{M})=0\).
Each graded summand therefore vanishes.
This argument requires neither an equivalence of ordinary module
categories nor compatibility between independently chosen gradings.

\section{The main theorem and applications}\label{sec:main}

\subsection{Proof of the main theorem and its consequences}

The following lemma is the Artin local case of
Avramov--Buchweitz's self-extension theorem.

\begin{lemma}[{\cite[Theorem 4.2]{AB00}}]\label{lem:commutative}
Let \(C=k[t_1,\ldots,t_c]/(t_1^{a_1},\ldots,t_c^{a_c})\), where
\(a_i\geq2\). A finite-dimensional left \(C\)-module \(L\) satisfying
\(\Ext_C^2(L,L)=0\) is free.
\end{lemma}

\begin{proof}
The ring \(C\simeq k[[t_1,\ldots,t_c]]/(t_1^{a_1},\ldots,t_c^{a_c})\)
is an Artin local complete intersection. The cited theorem gives
\(\pd_C L<\infty\). If \(L\ne0\), the Auslander--Buchsbaum formula yields
\(\pd_C L=\depth C-\depth_C L=0\). Thus \(L\) is projective, hence free
over the local ring \(C\). The zero module is free as well.
\end{proof}


\begin{theorem}\label{thm:main}
Let \(k\) be any field and let \(A=A(\mathbf a,\mathbf q)\) be as in
\eqref{eq:QCI}. For every finite-dimensional left \(A\)-module \(M\),
\[
 \Ext_A^1(M,M)=0=\Ext_A^2(M,M)
 \quad\Longrightarrow\quad M\text{ is projective}.
\]
\end{theorem}

\begin{proof}
We first assume that \(k\) is algebraically closed.

\smallskip\noindent
\emph{The intermediate algebras.}
For \(1\leq j\leq c\), put
\begin{equation*}
 q_{uv}^{(j)}=
 \begin{cases}q_{uv},&v\leq j,\\1,&v>j,\end{cases}
 \qquad
 B_j=\frac{k\langle x_1,\ldots,x_c\rangle}
 {(x_i^{a_i},\ x_ux_v-q_{uv}^{(j)}x_vx_u\ (u<v))}.
\end{equation*}
Here all power relations are included. Then \(B_c=A\), every
\(x_\ell\) with \(\ell>j\) is central in \(B_j\), and
\[
 B_1=C=k[x_1,\ldots,x_c]/(x_1^{a_1},\ldots,x_c^{a_c}).
\]
Each \(B_j\) has the basis of Lemma~\ref{lem:structure}.
For \(j\geq2\), use \(\deg_jx_i=\delta_{ij}\) and set
\begin{equation*}
 (x_i)\sigma_j=
 \begin{cases}q_{ij}^{-1}x_i,&i<j,\\x_i,&i\geq j.\end{cases}
\end{equation*}
This is a graded diagonal automorphism. Writing \(*_j\) for the
twisted multiplication, we have, for \(i<j\),
\begin{equation*}
 x_i*_jx_j=q_{ij}^{-1}x_ix_j=x_jx_i=x_j*_jx_i.
\end{equation*}
Two generators different from \(x_j\) both have degree zero, so their
products are unchanged. If \(\ell>j\), then \(\sigma_j\) fixes
\(x_\ell\), and \(x_\ell\) still commutes with \(x_j\).
All power relations are preserved, since a generator different from
\(x_j\) has degree zero and \(\sigma_j\) fixes \(x_j\).
Consequently there is a graded homomorphism
\begin{equation*}
 \phi_j:B_{j-1}\longrightarrow B_j^{\sigma_j},\qquad
 (x_i)\phi_j=x_i,
\end{equation*}
where both sides carry the \(j\)-th coordinate grading. On ordered
monomials it is given by
\begin{equation*}
 (x^\alpha)\phi_j=\lambda_j(\alpha)x^\alpha,
 \qquad \lambda_j(\alpha)=\prod_{i<j}q_{ij}^{-\alpha_i\alpha_j}.
\end{equation*}
Indeed, each of the \(\alpha_j\) occurrences of \(x_j\), when multiplied
in on the right, contributes \(\prod_{i<j}q_{ij}^{-\alpha_i}\).
All these coefficients are nonzero, so \(\phi_j\) maps the standard
basis to a basis and is an isomorphism.

\smallskip\noindent
\emph{The modules and the induction.}
Start with \(M_c=M\). Suppose \(M_j\in\umod{B_j}\) has been constructed,
with both \(\Ext_{B_j}^1(M_j,M_j)\) and \(\Ext_{B_j}^2(M_j,M_j)\) zero.
Theorem~\ref{thm:AO} provides a compatible \(j\)-th coordinate grading
\[
 \widetilde M_j=\bigoplus_{r\in\Z}(M_j)_r^{[j]}.
\]
Let \(R_j:\gmod{B_j^{\sigma_j}}\to\gmod{B_{j-1}}\) be restriction along
\(\phi_j\). Define \(M_{j-1}\) to be the ordinary module underlying
\(R_jF_{\sigma_j}(\widetilde M_j)\). For \(m\in(M_j)_r^{[j]}\),
\begin{equation*}
 x_i\cdot_{j-1}m=
 \begin{cases}
 q_{ij}^{-r}(x_i\cdot_jm),&i<j,\\
 x_i\cdot_jm,&i\geq j.
 \end{cases}
\end{equation*}
For an ordered monomial the full formula, including \(\phi_j\), is
\begin{equation*}
 x^\alpha\cdot_{j-1}m=
 \left(\prod_{i<j}q_{ij}^{-\alpha_i(r+\alpha_j)}\right)
 (x^\alpha\cdot_jm),
 \qquad m\in(M_j)_r^{[j]}.
\end{equation*}

Every power \(\sigma_j^d\) is diagonal. Proposition~\ref{prop:diagonal}
therefore gives \(\tw{\sigma_j^d}{M_j}\simeq M_j\) as ordinary modules
for all \(d\in\Z\). Proposition~\ref{prop:transport}, with
\(S=\{1,2\}\), and restriction along \(\phi_j\) now yield
\begin{gather}\label{eq:induction-invariants}
 \dim_kM_{j-1}=\dim_kM_j,\qquad
 \Ext_{B_{j-1}}^n(M_{j-1},M_{j-1})=0\quad(n=1,2),\\
 M_{j-1}\text{ is projective}
 \quad\Longleftrightarrow\quad M_j\text{ is projective}.\notag
\end{gather}
In particular, the new ordinary module \(M_{j-1}\) is rigid. After
forgetting its current grading, Theorem~\ref{thm:AO} can be applied
again for the \((j-1)\)-st coordinate grading. No compatibility between
successive choices is required. This proves that the induction is
legitimate at every stage.

After \(c-1\) steps, \(M_1\) is a \(C\)-module with
\(\Ext_C^2(M_1,M_1)=0\). Lemma~\ref{lem:commutative} makes \(M_1\)
free. Reflecting projectivity through \eqref{eq:induction-invariants}
gives projectivity of \(M_c=M\). If \(c=1\), there are no twists and
Lemma~\ref{lem:commutative} applies directly.

\smallskip\noindent
\emph{Descent to the original field.}
For arbitrary \(k\), choose an algebraic closure \(K\). The algebra
\(A_K=K\otimes_kA\) is again the quantum complete intersection
\eqref{eq:QCI} over \(K\). By Lemma~\ref{lem:base-change}, \(M_K\)
has zero first and second self-extension groups. The algebraically closed case shows that \(M_K\) is projective.
Lemma~\ref{lem:base-change} then implies that \(M\) is projective.
\end{proof}

\begin{thebibliography}{99}
\bibitem[AB00]{AB00} {\sc L.~L.~Avramov} and {\sc R.-O.~Buchweitz}, \href{https://doi.org/10.1007/s002220000090}{Support varieties and cohomology over complete intersections}, \emph{Invent. Math.} \textbf{142} (2000), 285--318.
\bibitem[AO13]{AO13} {\sc C.~Amiot} and {\sc S.~Oppermann}, \href{https://doi.org/10.1093/imrn/rns010}{The image of the derived category in the cluster category}, \emph{Int. Math. Res. Not. IMRN} \textbf{2013} (2013), no.~4, 733--760.
\bibitem[Hazrat16]{Hazrat16} {\sc R.~Hazrat}, \href{https://doi.org/10.1017/CBO9781316717134}{Graded Rings and Graded Grothendieck Groups}, \emph{London Mathematical Society Lecture Note Series, vol.~435, Cambridge University Press, Cambridge}, 2016.
\bibitem[Zhang96]{Zhang96} {\sc J.~J.~Zhang}, \href{https://doi.org/10.1112/plms/s3-72.2.281}{Twisted graded algebras and equivalences of graded categories}, \emph{Proc. London Math. Soc.} (3) \textbf{72} (1996), no.~2, 281--311.
\end{thebibliography}
\end{document}
